\chapter{Requirements and Analysis}
\label{chap:req_analysis}
This chapter defines the functional and non-functional requirements of the proposed system, establishing a framework for development and a benchmark for final evaluation. These requirements served as measurable milestones to track progress and ensured that the project aligned with its core objectives. To provide a structured implementation roadmap, each requirement was prioritised using the MoSCoW framework, alongside an estimation of the relative development effort required.

The analysis breaks the overarching problem into distinct, manageable deliverables, ranging from data aggregation and model training to voice interface design. Furthermore, this chapter addresses the legal and ethical considerations relevant to public datasets and esports data, ensuring the project adhered to established data protection standards and intellectual property rights.

\section{Project Requirements}
The following project requirements were defined and prioritised using the MoSCoW framework. This categorisation ensured that the core functionality of the Alexa skill was prioritised while secondary features that could enhance the user experience were still identified. Each requirement was further assigned a difficulty rating on a scale of 1--5, representing the estimated technical complexity and development time. Requirements classified as `Should' or `Could' represent non-essential enhancements that were planned for iterative implementation, provided that the primary objectives were met within the project timeline.

\subsection{Functional Requirements}
{
\renewcommand{\arraystretch}{1.3}
\begin{longtable}{@{} p{0.15\linewidth} p{0.65\linewidth} p{0.15\linewidth} @{} }
\caption{Functional Requirements prioritised by the MoSCoW method.} \label{tab:functional_req} \\
\toprule
\textbf{Priority} & \textbf{Requirement} & \textbf{Difficulty} \\
\midrule
\endfirsthead

\toprule
\textbf{Priority} & \textbf{Requirement} & \textbf{Difficulty} \\
\midrule
\endhead

\bottomrule
\endfoot

\cellcolor{green!15}\textbf{Must} & Request a prediction for an upcoming match between two recognised teams. & 2 \\
\cellcolor{green!15}\textbf{Must} & Receive a meaningful error message if teams or matches are unrecognised. & 2 \\
\cellcolor{green!15}\textbf{Must} & Exit the skill at any time via intuitive commands (e.g., ``Alexa, stop''). & 1 \\
\cellcolor{green!15}\textbf{Must} & Predict and announce the winning team. & 2 \\
\cellcolor{green!15}\textbf{Must} & Provide a confidence percentage for the prediction (e.g., ``65\% chance''). & 2 \\
\cellcolor{green!15}\textbf{Must} & Access a `help' intent providing a summary of available commands. & 2 \\
\cellcolor{green!15}\textbf{Must} & Support team names provided in any order within the utterance. & 2 \\ \addlinespace
\cellcolor{yellow!15}\textbf{Should} & Confirm team names if the initial voice input is ambiguous. & 3 \\
\cellcolor{yellow!15}\textbf{Should} & Provide reasoning or key metrics behind a specific prediction upon request. & 4 \\
\cellcolor{yellow!15}\textbf{Should} & Query historical results of previous matches between two teams. & 3 \\
\cellcolor{yellow!15}\textbf{Should} & Support multiple prediction requests within a single session. & 2 \\ \addlinespace
\cellcolor{orange!15}\textbf{Could} & Request secondary statistics such as head-to-head performance or win/loss ratios. & 3 \\
\end{longtable}
}

\subsection{Non-Functional Requirements}
{
\renewcommand{\arraystretch}{1.3}
\begin{longtable}{@{} p{0.15\linewidth} p{0.65\linewidth} p{0.15\linewidth} @{} }
\caption{Non-Functional Requirements prioritised by the MoSCoW method.} \label{tab:non_functional_req} \\
\toprule
\textbf{Priority} & \textbf{Requirement} & \textbf{Difficulty} \\
\midrule
\endfirsthead

\toprule
\textbf{Priority} & \textbf{Requirement} & \textbf{Difficulty} \\
\midrule
\endhead

\bottomrule
\endfoot

\cellcolor{green!15}\textbf{Must} & Intent recognition accuracy of at least 80\% across natural language utterances. & 3 \\
\cellcolor{green!15}\textbf{Must} & System response time (latency) must not exceed six seconds. & 4 \\
\cellcolor{green!15}\textbf{Must} & Machine learning model must achieve a minimum predictive accuracy of 60\%. & 5 \\ \addlinespace
\cellcolor{yellow!15}\textbf{Should} & Maintain a conversational tone and contextual awareness throughout the session. & 4 \\
\cellcolor{yellow!15}\textbf{Should} & Maintain session state to facilitate follow-up inquiries. & 4 \\
\cellcolor{yellow!15}\textbf{Should} & Robustness against minor phonetic variations in team name pronunciation. & 3 \\
\end{longtable}
}

\section{Evaluation Strategy}
The evaluation of this project was structured to verify that all functional and non-functional requirements had been satisfied while validating the system's efficacy as a predictive tool. Addressing these areas ensured that the final deliverable was both technically sound and capable of providing a meaningful and reliable experience for the end-user.

\subsection{Machine Learning Model Efficacy}
To evaluate the extent to which the project had satisfied the aim of providing accurate match predictions, the machine learning models were subjected to various statistical testing. The primary validation method involved partitioning the VCT dataset into a 70\% training set, a 20\% hold-out testing set, and a 10\% final validation set, allowing for an assessment of the model's ability to generalise to unseen data. Beyond simple accuracy, the project utilised a confusion matrix to derive precision, recall, and the F1-score. These metrics were vital for understanding the model's balance between correctly identifying winners and avoiding false positives, which is a more granular measure of success than a single percentage.

\subsection{Alexa Skill Testing}
The non-functional requirements regarding latency and reliability were critical for maintaining a seamless VUI. Verification was conducted using AWS CloudWatch to monitor the execution time of the Lambda function, tracking the duration from the moment the Alexa Voice Service identified a user's intent to the final delivery of the prediction. Success in this area was defined by maintaining an average response time of under six seconds, ensuring the system operated safely within Alexa's strict eight-second timeout limit. Furthermore, upon establishing an extensive collection of utterances, the coverage of the natural language understanding was verified by releasing the skill into a beta-testing environment. Potential users were identified and contacted in-line with the University's research ethics policies, and asked to try out the skill without knowledge of the existing utterances. This evaluated the system's ability to map real requests to the correct intent, targeting a minimum successful recognition rate of 80\%.

\subsection{User Experience}
The final layer of evaluation focused on qualitative satisfaction and the overall fulfilment of the project's original aims. A post-development user survey, conducted under the previously mentioned ethical guidelines, allowed a cohort of users to interact with the skill and provide feedback via Likert-scale questionnaires. This assessment covered the perceived naturalness of the conversation and the clarity of the logic provided by the model.

\section{Analysis}
This section details the implementation steps required to fulfil the requirements established in \autoref{tab:functional_req} and \autoref{tab:non_functional_req}. 

\subsection{Key Dataset Tables}
\label{sub:key-tables}
Analysis of the dataset revealed a subset of tables possessing high relational and predictive utility. These tables were categorised into two distinct groups based on their function: those serving as the primary source of predictive features, and those providing the identifier mappings essential for consolidating said feature-rich tables. 

In the absence of a unique ID, the combination of match descriptors (\texttt{Tournament, Stage, Match Type, Match Name, Map}) was not guaranteed to provide a unique reference. For example, three games played in a match between the teams `Acend' and `Envy' might resemble the following:

\begin{table}[H]
\centering
\small
\begin{tabular}{lllll}
\toprule
\textbf{Tournament} & \textbf{Stage} & \textbf{Match Type} & \textbf{Match Name} & \textbf{Map} \\
\midrule
Valorant Champions 2021 & Group Stage & Winner's (A) & Acend vs Envy & Icebox \\
Valorant Champions 2021 & Group Stage & Winner's (A) & Acend vs Envy & Haven \\
Valorant Champions 2021 & Group Stage & Winner's (A) & Acend vs Envy & Icebox \\
\bottomrule
\end{tabular}
\caption{Raw match data exhibiting non-unique map entries.}
\label{tab:raw_match_data}
\end{table}

In this series, the map `Icebox' was played twice, rendering the combination of attributes non-unique. However, the first key table, \texttt{combined\_matches\_games\_ids.csv}, provided the identifiers for each match and game, transforming the data into the following structure:

\begin{table}[H]
\centering
\small
\begin{tabular}{llll >{\columncolor{orange!30}}l l >{\columncolor{orange!30}}l}
\toprule
\textbf{Tournament} & \textbf{Stage} & \textbf{Match Type} & \textbf{Match Name} & \textbf{Match ID} & \textbf{Map} & \textbf{Game ID} \\
\midrule
Val. Champ. 21 & Group Stage & Winner's (A) & Acend vs Envy & 51282 & Icebox & 57948 \\
Val. Champ. 21 & Group Stage & Winner's (A) & Acend vs Envy & 51282 & Haven & 57949 \\
Val. Champ. 21 & Group Stage & Winner's (A) & Acend vs Envy & 51282 & Icebox & 57950 \\
\bottomrule
\end{tabular}
\caption{Refined data structure utilising \texttt{Match ID} and \texttt{Game ID}.}
\label{tab:refined_match_data}
\end{table}

The \texttt{combined\_matches\_games\_ids.csv} file acted as the primary relational bridge. It mapped match descriptors, such as \texttt{Match Name} and \texttt{Map}, to their unique IDs (e.g., \texttt{Match ID}, \texttt{Game ID}). This file was crucial for ensuring that performance data across the various sources were unified into a single, cohesive record for every match-map instance. Although the following key files did not natively contain the IDs of the matches or games, they did contain the aforementioned descriptors, which were used to locate the corresponding IDs in \texttt{combined\_matches\_games\_ids.csv}. They also provided the high-granularity performance data required for feature engineering:

\begin{itemize}
    \item \textbf{Match Scores (\texttt{map\_scores.csv}):} This served as the base table to build upon to produce the final feature set. It provided a row per match, alongside the primary outcome variables (scores and win/loss state) that the model aimed to predict.
    \item \textbf{Player Overview (\texttt{overview.csv}):} This file contained a granular player-level structure, offering advanced metrics such as \texttt{Average Combat Score} (ACS) and \texttt{KAST} percentage, as discussed in \autoref{tab:key-dataset-features}.
    \item \textbf{Economic Statistics (\texttt{eco\_stats.csv}):} This table provided round-level economic statistics, including the type of round (e.g., the first round of the game), and the outcome of the round. These statistics were contained within the features \texttt{Type} and \texttt{Won} respectively.
\end{itemize}

Before metrics from multiple tables could be consolidated, each source had to adhere to the same attributes for referencing by obtaining the IDs from \texttt{combined_matches_games_ids.csv}. The complete schemas for these tables can be found in \autoref{fig:tableSchema}.

\subsection{Data Pre-processing}
The foundation of the predictive engine was the Valorant Champions Tour dataset sourced from Kaggle \cite{ryanluong_dataset}, which, as established in \autoref{fig:dataset-folder-structure}, spanned professional match history from 2021 to 2025. A structural analysis of the source files revealed a consistent directory architecture across each annual folder, which facilitated the use of parameterised path references. By utilising variable strings (e.g., transitioning from \texttt{2021_ids} to \texttt{2025_ids}), a programmatic pipeline was established to ingest multi-year data without manual intervention.

To ensure relational integrity across the disparate tables, the project utilised the master ID system already present. Each entity type---including players, teams, and matches---was assigned a unique identifier that served as a primary key. Before feature derivation could commence, these master tables were verified for completeness; any missing identifiers would result in reference errors during the join operations.

The data preparation workflow utilised the \textit{pandas} library \cite{pandas_site} to perform these relational operations. Data from individual years was read into DataFrame objects and standardised to ensure consistent schema mapping. These DataFrames were subsequently merged using the \texttt{concat()} function. A critical consideration during this merge was the mitigation of data redundancy; because professional players and teams often persist across multiple seasons, the union of annual tables inevitably introduced duplicate records. These were identified and purged using the \texttt{drop\_duplicates()} function, ensuring each primary key remained unique.

Furthermore, the dataset was audited for data sparsity, often represented as `NaN' (Not a Number) or null values, which can significantly degrade the performance of machine learning algorithms \cite{11226584}. To quantify missing data, a diagnostic check using the \texttt{isnull()} and \texttt{sum()} functions was carried out. Where data was missing, an imputation strategy was implemented: by replacing these null gaps with feature averages, the final training set was ensured to be mathematically sound.

\subsection{Feature engineering}
Once the dataset had been ingested and validated for integrity, the project moved into the feature engineering phase. In this phase, raw match statistics were transformed into a format that captured the ``form'' and ``momentum'' of the competing teams. To achieve this, a rolling average strategy was implemented. Rather than providing the model with a static snapshot of a team's history, the system calculated performance metrics (such as Win Rate, ACS, and KAST) over a sliding window of the previous $n$ matches, ensuring that the model's predictions were based on recent competitive standing rather than outdated historical data, reflecting the continuous evolution of professional strategies in Valorant.

The transformation process also involved the calculation of ``Differential Features.'' Machine learning models, particularly tree-based ensembles, often benefit from seeing the relative strength between two competitors. Therefore, the implementation involved calculating the difference between Team A and Team B for key metrics. This reduced the dimensionality of the input while highlighting the competitive gap between the two teams, which was a primary indicator of match outcomes.

\subsection{Model Development}
With a structured table of rolling features and differentials established, the project transitioned to the model engineering phase. This involved defining the architecture for the chosen supervised learning algorithms: Logistic Regression and Gradient Boosted Trees. The implementation followed a train, test, tweak cycle, where initially, a baseline was established using Logistic Regression to provide a point of comparison for the more complex ensemble model.

The training process incorporated hyperparameter tuning to optimise model performance. Parameters such as the number of estimators, maximum depth, and learning rate were adjusted manually. Following each training iteration, the model was evaluated against the 20\% hold-out test set---using the metrics defined in the evaluation strategy---to determine the strength of the current feature set and the generalisability of the model. This iterative process continued until the 60\% accuracy threshold was consistently met or exceeded, upon which the model was evaluated against the final 10\% validation set. Once an optimal model was produced, it was exported using the \textit{m2cgen} library to allow for efficient loading within the AWS Lambda environment.

\section{Legal and Ethical Considerations}
The development and deployment of the predictive esports skill involved several legal and ethical dimensions, ranging from data ownership to the social implications of ML-driven forecasting.

\subsection{Data Licensing and Intellectual Property}
The primary dataset used for model training was sourced from Kaggle \cite{ryanluong_dataset} and was utilised under its specified MIT licence. Furthermore, match data referenced from VLR.gg \cite{vlr_gg} was used strictly for non-commercial, academic research purposes. The project acknowledged the intellectual property of Riot Games as the developer of Valorant; however, the use of public match statistics for algorithmic analysis falls well within the provisions of fair dealing for non-commercial research.

\subsection{Data Protection}
As an Alexa-hosted application, the skill adhered to the ``Privacy by Design'' principle. The system did not record or store personally identifiable information. Voice recordings were processed by Amazon's Alexa Voice Service (AVS), and the skill only received the resulting intent and slot data. No persistent user data was stored on AWS Lambda beyond the duration of the active session, ensuring compliance with the GDPR and the UK Data Protection Act 2018.

\subsection{Gambling Ethics}
A fundamental ethical requirement of this project was to manage user expectations regarding the fallibility of machine learning. It was imperative to communicate that the model's outputs were probabilistic estimates based on historical data, rather than absolute certainties. To mitigate the risk of the skill being misused as a financial advisory tool for gambling, a clear disclaimer was included on the skill's description page. 

Furthermore, it was acknowledged that the model might exhibit inherent algorithmic bias toward historically dominant teams due to the higher volume of associated training data. To ensure transparency, the skill was designed to provide a confidence percentage (e.g., ``65\% chance'') to highlight the uncertainty inherent in the prediction. This design choice encouraged users to view the system as an informational tool rather than a definitive source of truth.

\subsection{Content Sensitivity}
While Valorant is rated PEGI 16 due to its depiction of tactical combat and firearm usage, Alexa-enabled devices are frequently positioned in communal household environments accessible to younger children. Ethically, it was imperative that the skill remained accessible without exposing unintended audiences to mature themes.

To maintain a neutral and inclusive user experience, the skill's response logic purposefully avoided violent terminology. References to ``kills'' or ``deaths'' were replaced with neutral descriptors such as ``eliminations,'' ``round wins,'' or ``match outcomes.'' By focusing on the strategic and competitive results rather than the violent mechanics of the game, the skill adhered to the safety standards expected of a family-oriented smart speaker ecosystem.

\subsection{Research Ethics and Participant Consent}
The final phase of this project involved a post-development user survey to evaluate the system's efficacy and user experience. As such, formal ethics approval had been obtained from the University Research Ethics Committee prior to any data collection. The study was designed to exclude sensitive topics and vulnerable participants. All data gathered was anonymised, and no personal identifiers were stored, ensuring the research was conducted in accordance with the university's ethical framework for human participation. The approved ethics applications for the use of the dataset and the execution of the user feedback questionnaire can be found in \autoref{app:dataset_application} and \autoref{app:questionnaire_application}, respectively.